\documentclass[10pt,a4paper]{amsart}
\usepackage{amsmath,amssymb,amsthm,graphicx}
% Retained minimal preamble for the Stacks Project chapter sites-cohomology.tex (Cohomology on Sites).
% Verbatim excerpt of the chapter preamble: Xy-pic setup (native preamble lines 14-18) and the native
% theorem declarations with the mathematical macros (native preamble lines 154-217). Copied mechanically;
% nothing added, removed or redefined. The wrapper class and the amsmath/amssymb/amsthm/graphicx base
% packages are supplied by the assembler, so the class/externaldocument part of the native preamble is omitted.
% For commutative diagrams we use Xy-pic
\usepackage[all]{xy}

% We use 2cell for 2-commutative diagrams.
\xyoption{2cell}
% Theorem environments.
%
\theoremstyle{plain}
\newtheorem{theorem}[subsection]{Theorem}
\newtheorem{proposition}[subsection]{Proposition}
\newtheorem{lemma}[subsection]{Lemma}

\theoremstyle{definition}
\newtheorem{definition}[subsection]{Definition}
\newtheorem{example}[subsection]{Example}
\newtheorem{exercise}[subsection]{Exercise}
\newtheorem{situation}[subsection]{Situation}

\theoremstyle{remark}
\newtheorem{remark}[subsection]{Remark}
\newtheorem{remarks}[subsection]{Remarks}

\numberwithin{equation}{subsection}

% Macros
%
\def\lim{\mathop{\mathrm{lim}}\nolimits}
\def\colim{\mathop{\mathrm{colim}}\nolimits}
\def\Spec{\mathop{\mathrm{Spec}}}
\def\Hom{\mathop{\mathrm{Hom}}\nolimits}
\def\Ext{\mathop{\mathrm{Ext}}\nolimits}
\def\SheafHom{\mathop{\mathcal{H}\!\mathit{om}}\nolimits}
\def\SheafExt{\mathop{\mathcal{E}\!\mathit{xt}}\nolimits}
\def\Sch{\mathit{Sch}}
\def\Mor{\mathop{\mathrm{Mor}}\nolimits}
\def\Ob{\mathop{\mathrm{Ob}}\nolimits}
\def\Sh{\mathop{\mathit{Sh}}\nolimits}
\def\NL{\mathop{N\!L}\nolimits}
\def\CH{\mathop{\mathrm{CH}}\nolimits}
\def\proetale{{pro\text{-}\acute{e}tale}}
\def\etale{{\acute{e}tale}}
\def\QCoh{\mathit{QCoh}}
\def\Ker{\mathop{\mathrm{Ker}}}
\def\Im{\mathop{\mathrm{Im}}}
\def\Coker{\mathop{\mathrm{Coker}}}
\def\Coim{\mathop{\mathrm{Coim}}}

% Boxtimes
%
\DeclareMathSymbol{\boxtimes}{\mathbin}{AMSa}{"02}

%
% Macros for moduli stacks/spaces
%
\def\QCohstack{\mathcal{QC}\!\mathit{oh}}
\def\Cohstack{\mathcal{C}\!\mathit{oh}}
\def\Spacesstack{\mathcal{S}\!\mathit{paces}}
\def\Quotfunctor{\mathrm{Quot}}
\def\Hilbfunctor{\mathrm{Hilb}}
\def\Curvesstack{\mathcal{C}\!\mathit{urves}}
\def\Polarizedstack{\mathcal{P}\!\mathit{olarized}}
\def\Complexesstack{\mathcal{C}\!\mathit{omplexes}}
% \Pic is the operator that assigns to X its picard group, usage \Pic(X)
% \Picardstack_{X/B} denotes the Picard stack of X over B
% \Picardfunctor_{X/B} denotes the Picard functor of X over B
\def\Pic{\mathop{\mathrm{Pic}}\nolimits}
\def\Picardstack{\mathcal{P}\!\mathit{ic}}
\def\Picardfunctor{\mathrm{Pic}}
\def\Deformationcategory{\mathcal{D}\!\mathit{ef}}
\usepackage{cleveref}
\crefname{originallemma}{original lemma}{original lemmas}
\Crefname{originallemma}{Original lemma}{Original lemmas}
\makeatletter
\newlabel{lemma-h1-extensions}{{03F1}{source}{}{sourcepointer-1}{}}
\newlabel{lemma-h1-extensions@cref}{{[originallemma][0][]03F1}{[1][0][]source}{}{}{}}
\makeatother
\crefname{originallemma}{original lemma}{original lemmas}
\Crefname{originallemma}{Original lemma}{Original lemmas}
\makeatletter
\newlabel{sites-modules-lemma-flat-tor-zero}{{03EX}{source}{}{sourcepointer-2}{}}
\newlabel{sites-modules-lemma-flat-tor-zero@cref}{{[originallemma][0][]03EX}{[1][0][]source}{}{}{}}
\makeatother
\title{The first cohomology of an O-module agrees with the first cohomology of its underlying abelian sheaf}
\author{The Stacks Project Authors}
\date{stacks-project:sites-cohomology.tex:local-snapshot-20261011; GNU FDL 1.2 or later}
\begin{document}
\maketitle
\noindent\textit{Editorial note.} All delivered mathematical text is the author's own, copied from the pinned Stacks Project chapter file sites-cohomology.tex; no statement or proof step was rewritten, reordered or removed. The author writes 'The converse is a little more fun' and omits the double-composition check; both are reproduced as written.\par\medskip
\noindent\textit{Editorial note.} The preamble is the chapter's own retained minimal preamble (sources/stacks-sites-cohomology/retained-preamble.tex), a verbatim excerpt of the native Xy-pic setup and of the native theorem declarations and mathematical macros; the Xy-pic xy environment is required because the counted proof displays two maps of short exact sequences. No macro was invented or redefined.\par\medskip
\noindent\textit{Editorial note.} The same-chapter reference is located at its \textbackslash{}label line inside the pinned chapter file; the cross-chapter pointer is located in the project's own tag list staged in the snapshot (sources/stacks-chow/tags.txt), at the verbatim line carrying its canonical tag.\par\medskip
\noindent\textit{Allowed standard prerequisites (editorial).}\begin{itemize}
\item The bijection between Ext\textasciicircum{}1 and H\textasciicircum{}1 of the structure sheaf. Same chapter, Lemma labelled lemma-h1-extensions (source line 372, Stacks tag 03F1); invoked at proof line 442. Supplies the translation of both sides of the asserted isomorphism into groups of extensions, which is the first step of the proof; the two comparison constructions are the counted content.
\item Tensoring with a flat module preserves exactness. Modules on Sites, Lemma labelled sites-modules-lemma-flat-tor-zero (canonical Stacks tag 03EX); invoked at proof lines 480-481. Supplies exactness of the sequence obtained from an extension of Z by F\_ab after tensoring with O; the pushout construction on top of it is retained in the counted span.
\end{itemize}\medskip
\ifdefined\hypertarget\hypertarget{sourcepointer-1}{}\fi
\noindent\textit{Source reference.} Original lemma 03F1: sources/stacks-sites-cohomology/pinned-original-source.tex, source lines 372--372. Same-chapter pointer invoked at proof line 442 to translate both H\textasciicircum{}1 groups into extension groups; the two comparison constructions that form the counted argument are retained verbatim in the delivered body.\par\medskip
\ifdefined\hypertarget\hypertarget{sourcepointer-2}{}\fi
\noindent\textit{Source reference.} Original lemma 03EX: sources/stacks-chow/tags.txt, source lines 4472--4472. Cross-chapter standard prerequisite invoked once at proof lines 480-481 for exactness after tensoring with O; the pushout construction on top of it is retained in the counted span.\par\medskip
\begin{lemma}
\label{lemma-h1-mod-ab-agree}
Let $(\mathcal{C}, \mathcal{O})$ be a ringed site.
Let $\mathcal{F}$ be a sheaf of $\mathcal{O}$-modules on $\mathcal{C}$.
Let $\mathcal{F}_{ab}$ denote the underlying sheaf of abelian
groups. Then there is a functorial isomorphism
$$
H^1(\mathcal{C}, \mathcal{F}_{ab})
=
H^1(\mathcal{C}, \mathcal{F})
$$
where the left hand side is cohomology computed in
$\textit{Ab}(\mathcal{C})$ and the right hand side
is cohomology computed in $\textit{Mod}(\mathcal{O})$.
\end{lemma}

\begin{proof}
Let $\underline{\mathbf{Z}}$ denote the constant sheaf
$\mathbf{Z}$. As
$\textit{Ab}(\mathcal{C}) = \textit{Mod}(\underline{\mathbf{Z}})$
we may apply
Lemma \ref{lemma-h1-extensions}
twice, and it follows that we have to show
$$
\Ext^1_{\textit{Mod}(\mathcal{O})}(\mathcal{O}, \mathcal{F})
=
\Ext^1_{\textit{Mod}(\underline{\mathbf{Z}})}(
\underline{\mathbf{Z}}, \mathcal{F}_{ab}).
$$
Suppose that $0 \to \mathcal{F} \to \mathcal{E} \to \mathcal{O} \to 0$
is an extension in $\textit{Mod}(\mathcal{O})$. Then we can use
the obvious map of abelian sheaves
$1 : \underline{\mathbf{Z}} \to \mathcal{O}$
and pullback to obtain an extension $\mathcal{E}_{ab}$, like so:
$$
\xymatrix{
0 \ar[r] &
\mathcal{F}_{ab} \ar[r] \ar@{=}[d] &
\mathcal{E}_{ab} \ar[r] \ar[d] &
\underline{\mathbf{Z}} \ar[r] \ar[d]^{1} &
0 \\
0 \ar[r] &
\mathcal{F} \ar[r] &
\mathcal{E} \ar[r] &
\mathcal{O} \ar[r] &
0
}
$$
The converse is a little more fun. Suppose that
$0 \to \mathcal{F}_{ab} \to \mathcal{E}_{ab} \to \underline{\mathbf{Z}} \to 0$
is an extension in $\textit{Mod}(\underline{\mathbf{Z}})$.
Since $\underline{\mathbf{Z}}$ is a flat $\underline{\mathbf{Z}}$-module
we see that the sequence
$$
0 \to \mathcal{F}_{ab} \otimes_{\underline{\mathbf{Z}}} \mathcal{O}
\to \mathcal{E}_{ab} \otimes_{\underline{\mathbf{Z}}} \mathcal{O}
\to \underline{\mathbf{Z}} \otimes_{\underline{\mathbf{Z}}} \mathcal{O}
\to 0
$$
is exact, see
Modules on Sites, Lemma \ref{sites-modules-lemma-flat-tor-zero}.
Of course
$\underline{\mathbf{Z}} \otimes_{\underline{\mathbf{Z}}} \mathcal{O}
= \mathcal{O}$.
Hence we can form the pushout via the ($\mathcal{O}$-linear) multiplication map
$\mu : \mathcal{F} \otimes_{\underline{\mathbf{Z}}} \mathcal{O}
\to \mathcal{F}$ to get an extension of $\mathcal{O}$ by
$\mathcal{F}$, like this
$$
\xymatrix{
0 \ar[r] &
\mathcal{F}_{ab} \otimes_{\underline{\mathbf{Z}}} \mathcal{O}
\ar[r] \ar[d]^\mu &
\mathcal{E}_{ab} \otimes_{\underline{\mathbf{Z}}} \mathcal{O}
\ar[r] \ar[d] &
\mathcal{O} \ar[r] \ar@{=}[d] &
0 \\
0 \ar[r] &
\mathcal{F} \ar[r] &
\mathcal{E} \ar[r] &
\mathcal{O} \ar[r] &
0
}
$$
which is the desired extension. We omit the verification that these
constructions are mutually inverse.
\end{proof}

\end{document}
